\documentclass[10pt, t, fleqn, aspectratio=43, %handout
]{beamer} %aspectratio can be 43 or 169
%\usepackage[english]{babel}
%\documentclass[10pt, t, fleqn, aspectratio=169]{beamer}
\usepackage[43]{SUslide}
\usepackage{graphicx}
\usepackage[skins,theorems]{tcolorbox}
\usepackage{tikz}
\usepackage{venndiagram}
\usepackage{ifthen}
\usepackage[all]{xy}
%\usetikzlibrary{positioning,shadows,backgrounds}
%\usepackage{empheq}
%\usepackage{animate}
%\usepackage{table}
%-------------------------------------------------------
% INFORMATION IN THE TITLE PAGE
%-------------------------------------------------------

\title[Short title] % [] is optional - is placed on the bottom of the sidebar on every slide
{ % is placed on the title page
      \textbf{Title}
}

\subtitle[]
{
      \textbf{Extra info}
}

\author[Short Author]
{      Author \\
      \alert{\ttfamily xxx@math.su.se}
}

\institute[]
{
10

      Department of Mathematics, Faculty of Science\\
      Stockholm University\\
  
  %there must be an empty line above this line - otherwise some unwanted space is added between the university and the country (I do not know why;( )
}

\date{\today}
\setbeamertemplate{caption}{\insertcaption}
%-------------------------------------------------------
% THE BODY OF THE PRESENTATION
%-------------------------------------------------------
\makeatletter
\newcommand\binomialCoefficient[2]{%
    % Store values 
    \c@pgf@counta=#1% n
    \c@pgf@countb=#2% k
    %
    % Take advantage of symmetry if k > n - k
    \c@pgf@countc=\c@pgf@counta%
    \advance\c@pgf@countc by-\c@pgf@countb%
    \ifnum\c@pgf@countb>\c@pgf@countc%
        \c@pgf@countb=\c@pgf@countc%
    \fi%
    %
    % Recursively compute the coefficients
    \c@pgf@countc=1% will hold the result
    \c@pgf@countd=0% counter
    \pgfmathloop% c -> c*(n-i)/(i+1) for i=0,...,k-1
        \ifnum\c@pgf@countd<\c@pgf@countb%
        \multiply\c@pgf@countc by\c@pgf@counta%
        \advance\c@pgf@counta by-1%
        \advance\c@pgf@countd by1%
        \divide\c@pgf@countc by\c@pgf@countd%
    \repeatpgfmathloop%
    \the\c@pgf@countc%
}
\makeatother
\newcommand\Stirling[2]{%
    % Store values 
    \c@pgf@counta=#1% n
    \c@pgf@countb=#2% k
    %
    
    % Recursively compute the coefficients
    \c@pgf@countc=1% will hold the result
    \c@pgf@countd=0% counter
    \pgfmathloop% c -> c*(n-i)/(i+1) for i=0,...,k-1
        \ifnum\c@pgf@countd<\c@pgf@countb%
        \multiply\c@pgf@countc by\c@pgf@counta%
        \advance\c@pgf@counta by-1%
        \advance\c@pgf@countd by1%
        \divide\c@pgf@countc by\c@pgf@countd%
    \repeatpgfmathloop%
    \the\c@pgf@countc%
}
\makeatother
\begin{document}



%\begin{frame}{Content}{}
%\tableofcontents
%\end{frame}

%-------------------------------------------------------
\section{A section}
%-------------------------------------------------------

\begin{frame}{A Slide}
%-------------------------------------------------------

  \begin{itemize}
    \item<1-> \textbf{Bok:} \emph{Discrete Mathematics, second edition} av Norman L. Biggs. Oxford University Press.
    \item<2-> \textbf{Tid och Plats:} Tisdagar och fredagar kl 12:45-14:30 Sal 15 Hus 5. Det ska ocks\aa\; vara \textbf{\"ovningar} p\aa\; tisdagar och fredagar kl 14:45-16:30.
    \item<3-> Det finns 2 stycken problemsamlingar (Webwork), med vardera 14 uppgifter. Avklaras alla uppgifter inom en problemsamling  ges 1,5 bonuspo\"ang p\aa\; tentamen.
  \end{itemize}
\end{frame}
\begin{frame}{Another slide}
%-------------------------------------------------------
\begin{definition}
A definition
\end{definition}
\begin{block}
{A block}
\end{block}
\begin{exampleblock}
{An example block}example
\end{exampleblock}
\begin{alertblock}
{An alert block}
\end{alertblock}
\end{frame}


%%%%%%%%%%%%%%%%%%%%%%%%%%%%%%%%%%%%%%%%%%%%%%%%%%%%%%%%%%%%%%%%%%%%%%%%%%%%%%%%%%%%%%%%%%
\section{Talsystem, \"Andliga och O\"andliga M\"angder}
%%%%%%%%%%%%%%%%%%%%%%%%%%%%%%%%%%%%%%%%%%%%%%%%%%%%%%%%%%%%%%%%%%%%%%%%%%%%%%%%%%%%%
\begin{frame}{Talsystem}
    An box\\
    \hfill\\
    \begin{tcolorbox}[colback=migrey!50,colframe=miblue]
     \hfill\\
 $$a_n\ldots a_1a_0\quad\text{betyder}\quad a_n\cdot 10^n+\cdots+a_1\cdot 10+a_0\cdot 10^0$$
\end{tcolorbox}

 
\end{frame}


\end{document}

%%% Local Variables:
%%% mode: latex
%%% TeX-master: t
%%% End:
